\begingroup
\fancyhead[L]{\small REMARKS}
\renewcommand{\headrulewidth}{0pt}
\par\noindent\hspace*{\parindent}On [III.]--[IV.]: The calculations for the special perturbations are found
mostly on sheets and to a lesser extent in notebooks. The entire text
and some of the numbers had to be supplied. The following material was taken from the originals.
On pp. 473--474 are the formulas found on the sheets belonging to the second
\par\smallskip\noindent\hfill 76\textsuperscript{*}\par
\noindent\SourceLeaf{610}calculation of the special perturbations.
On pp. 475--482 are, with a few exceptions, the numbers scattered partly on sheets and partly in
two notebooks (B\,c and B\,d). Some of these numbers were supplemented.
The numbers for the second calculation, pp. 483--488, are found without accompanying text
but in sequence, in a bundle of sheets.

\par\noindent\hspace*{\parindent}On [V.]: Notes [1.]--[3.] and [5.] had to be freely reconstructed throughout,
since the posthumous papers contain no account of the method of calculation.
They contain only the sheets with the calculations, and it was only after deciphering them that
the method could be made clear. In this the editor drew on
some preliminary work by \textsc{Schering} and \textsc{Schur}, who had already deciphered several sheets
and indicated the meaning of some of the numbers. Some of
\textsc{Schering}'s notes also suggested that he might have
understood most of the calculations. Unfortunately, those
notes were too incomplete to use in preparing the edition.
The worked example [4.], which appears in a notebook and is printed with only minor
changes, served as a guide. The presentation is arranged
to follow the »Theoria interpolationis«, Vol. III, p. 265, since
this was probably the approach \textsc{Gauss} took at first.
The interesting investigation in the »Quatrième Section« of the »Exposition«, p. 465,
therefore does not come into play here. Its connection with the numerical
calculation is plain, however, when the integrations for determining the
coefficients of the series are performed by mechanical quadrature.

\par\noindent\hspace*{\parindent}Of the numerical values in note [6.], those on pp. 507--508
were found after a long search in a small notebook, Schedae (A\,n). Those on pp. 509--513
are on separate sheets. They had to be supplemented in several places. In the tables
for \(T\), \(V\), \(W\), pp. 509--511, some values are missing and have been left unsupplied.
The calculation of the perturbations themselves, notes [7.]--[12.], and the comparison with
the oppositions, note [13.], are on separate sheets. Only as much was
printed as appeared necessary for an account without gaps.

\par\noindent\hspace*{\parindent}One of the sheets containing the integration of the perturbations of the epoch
\(\varepsilon\) has the following short calculation. The first of the
numbers given is the mean motion of Pallas:
\begingroup
\newcommand{\BNoteJ}{\text{\mdseries\AstroSymbols ♃}}
\newcommand{\BNoteP}{\text{\mdseries\AstroSymbols ⚴}}
\[
\begin{gathered}
\mathit{769\overset{\prime\prime}{,}202079}\\[3pt]
\textit{7 times this}=\mathit{5384\overset{\prime\prime}{,}414553}\\[3pt]
\underline{\mathit{18}\ \textit{m.\,m.}\ \BNoteJ=\mathit{5384{,}392272}}\\[3pt]
\mathit{18}\BNoteJ-\mathit{7}\BNoteP=-\mathit{0{,}022281}\\[3pt]
\BNoteJ\ \textit{sid. in one day}\ \mathit{299\overset{\prime\prime}{,}132904}.
\end{gathered}
\]
\endgroup
\par\noindent This seems to be the only indication of the first step in
\textsc{Gauss}'s discovery of the rational ratio of the two orbital periods. All that
can be inferred from it is that he had found the quantity \(18n'-7n\)
to be extremely small.

\par\noindent\hspace*{\parindent}For some of the underlying elements of Pallas, p. 508,
\textsc{Gauss} adopted slightly different values in the course of the calculation.
For example, in the integrations he used \(\tilde\omega-\AscendingNode=308^0\ 29'\ 23''\).
A few minor calculation errors have been corrected, but the calculation was not checked
because of the enormous amount of work involved.
\textsc{Gauss} seems to have carried out this first calculation of the general perturbations with less
care.
\clearpage
\endgroup
